% Generated draft: requires mathematical and source review.
\documentclass[11pt]{article}
\usepackage[margin=25mm]{geometry}
\usepackage{fontspec}
\setmainfont{Latin Modern Roman}
\usepackage{amsmath,amssymb,mathtools,unicode-math}
\setmathfont{Latin Modern Math}
\usepackage{enumitem,needspace,xurl,hyperref}
\setlength{\parindent}{0pt}
\setlength{\parskip}{5pt}
\setlength{\emergencystretch}{3em}
\newcommand{\N}{\mathbb{N}}
\newcommand{\Z}{\mathbb{Z}}
\newcommand{\Q}{\mathbb{Q}}
\newcommand{\R}{\mathbb{R}}
\newcommand{\C}{\mathbb{C}}
\newcommand{\F}{\mathbb{F}}
\newcommand{\A}{\mathbb{A}}
\newcommand{\PP}{\mathbb{P}}
\newcommand{\E}{\mathbb{E}}
\newcommand{\eps}{\varepsilon}
\newcommand{\dd}{\,\mathrm{d}}
\begin{document}
\section*{UnsolvedMath ID: 20002704}
\textbf{Source problem code:} catalogue entry
\subsection{Definitions and mathematical statement}
The aid identifies this extraction as two genuine remarks/questions following AIM Problem 1.3, not an independent Problem 5.3. The leading ``Problem 5.3, in free probability'' is an internal reference from surrounding material. The first question asks for a natural generalized nonnesting poset $NN_W^{(k)}$: in type $A$, it suggests $k$-divisible nonnesting set partitions under refinement, perhaps matching an order on co-filtered root-poset multichains in general. The second names Christos Athanasiadis and Stavros Garoufallidis as suggesting a $q$-version of Catalan combinatorics, but ends with an incomplete ``See'' and no reference or definition.

For the aid's concrete test let $k\geq1$. A set partition is $k$-divisible if each block size is a multiple of $k$; nonnesting means that its standard arcs joining consecutive elements of each block contain no nested pair $a<c<d<b$ with arcs $(a,b)$ and $(c,d)$. Order partitions by refinement. In rank one, consider these partitions of $[2k]$. They consist of the one-block maximum and exactly $k$ incomparable two-block partitions, so their poset is the $k$-fork $F_k$, with $k$ minimal elements and one maximum.

In contrast, the positive-root poset of $A_1$ is $\{\alpha\}$. Its geometric decreasing $k$-filter chains have the form
\[
F_1=\cdots=F_j=\{\alpha\},\qquad
F_{j+1}=\cdots=F_k=\varnothing\qquad(0\leq j\leq k).
\]
Root-addition constraints are vacuous here. Componentwise inclusion makes these $k+1$ choices a chain $C_{k+1}$. Therefore for every $k\geq2$,
\[
C_{k+1}\not\cong F_k,
\]
even though both have $k+1$ elements. This refutes only the immediate componentwise-inclusion order-isomorphism proposal, not every possible order on geometric multichains or every natural generalized nonnesting construction. For $k=1$ the rank-one obstruction disappears.

For $W=A_1^r$, each of the $r$ independent positive roots has a threshold $j_\ell\in\{0,\ldots,k\}$, so the same filter order is $C_{k+1}^r$ with natural rank $\rho=\sum_\ell j_\ell$. The aid's $q$-Fuss-Catalan normalization gives
\[
\sum_{F}q^{2\rho(F)}=(1+q^2+\cdots+q^{2k})^r
=\left(\frac{[2k+2]_q}{[2]_q}\right)^r,
\qquad [m]_q=1+q+\cdots+q^{m-1}.
\]
Thus twice the natural rank realizes this polynomial in the $A_1^r$ test family. It is not a general statistic theorem for all root systems, and it does not identify the incomplete source's exact intended $q$-construction.
\subsection{Short English statement}
The fragment asks for a generalized nonnesting order and mentions an undefined $q$-Catalan proposal. The aid shows a rank-one obstruction: componentwise filter chains form a chain, while divisible nonnesting refinement forms a fork. A product of rank-one systems realizes a particular $q$-Fuss-Catalan polynomial by twice its filter-chain rank.
\subsection{Sources}
\begin{itemize}
\item Local UnsolvedMath registry, permanent ID $20002704$, \texttt{catalogue entry}; \texttt{aim-probability-notes.json}, \texttt{notes[145]}, extracted source item 5.3.
\item American Institute of Mathematics, \emph{Braid groups, clusters and free probability}, \url{https://aimath.org/WWN/braidgroups/braidgroups.pdf}. The aid locates the two substantive remarks after Problem 1.3 rather than at a standalone Problem 5.3.
\item Original attribution: Christos Athanasiadis and Stavros Garoufallidis for a $q$-version of Catalan combinatorics. The citation following ``See'' is missing. No authoritative separate reference for the rank-one obstruction or product enumerator is supplied.
\item Research attempt 1: partial result, problem partially solved at run, rigor label ``proved,'' novelty confidence low. Registry classification partial progress, status partially solved, statement status exact, despite the contextual extraction and truncated reference.
\end{itemize}
\end{document}
